\subsubsection{External Interface Requirements}

\resetreqcounter{GeIR}
\begin{srstable}[
  caption = {Требования к общим интерфейсам.},
]{colspec = {Q[c, m] X[j, m]}}
  Идентификатор & Описание \\
  \reqlabel[ref:gei1]{GeIR} & Система \textit{Должна} реализовывать поисковую функциональность посредством интеграции с программируемой поисковой системой Google. \\
  \reqlabel[ref:gei2]{GeIR} & Система \textit{Должна} обеспечивать отправку поискового запроса из поля поиска, расположенного в шапке Сайта, к индексу Google, ограниченному доменом \texttt{magya-online.ru}. \\
  \reqlabel[ref:gei3]{GeIR} & Система \textit{Должна} отображать результаты поиска на отдельной странице в стандартном формате, предоставляемом поисковой системой Google. \\
\end{srstable}

\resetreqcounter{CmIR}
\begin{srstable}[
  caption = {Требования к интерфейсу комментариев.},
]{colspec = {Q[c, m] X[j, m]}}
  Идентификатор & Описание \\
  \reqlabel[ref:cmi1]{CmIR} & Система \textit{Должна} осуществлять отправку данных формы комментария с использованием метода HTTP POST. \\
  \reqlabel[ref:cmi2]{CmIR} & Система \textit{Должна} принимать текстовые поля формы комментария без ограничения максимальной длины. \\
  \reqlabel[ref:cmi3]{CmIR} & Система \textit{Должна} устанавливать обязательность заполнения всех полей формы отправки комментария. \\
  \reqlabel[ref:cmi4]{CmIR} & Система \textit{Должна} отображать сообщение «Ваш комментарий добавлен. После рассмотрения модератором он будет размещён на сайте» при успешной обработке отправленной формы комментария. \\
  \reqlabel[ref:cmi5]{CmIR} & Система \textit{Должна} отображать сообщение «Введите все данные!» при обнаружении незаполненных полей в форме отправки комментария. \\
  \reqlabel[ref:cmi6]{CmIR} & Система \textit{Должна} предоставлять пагинацию комментариев в виде элемента управления с наименованием «Смотреть все комментарии». \\
\end{srstable}

\resetreqcounter{RtIR}
\begin{srstable}[
  caption = {Требования к интерфейсу рейтинга.},
]{colspec = {Q[c, m] X[j, m]}}
  Идентификатор & Описание \\
  \reqlabel[ref:rti1]{RtIR} & Система \textit{Должна} осуществлять отправку выбранного пользователем значения оценки незамедлительно после совершения взаимодействия с виджетом, без дополнительного подтверждения. \\
  \reqlabel[ref:rti2]{RtIR} & Система \textit{Должна} отображать сообщение «Ваш голос засчитан. Спасибо!» при успешной обработке отправленной оценки. \\
  \reqlabel[ref:rti3]{RtIR} & Система \textit{Должна} отображать среднюю арифметическую оценку и общее количество голосов, поданных пользователями. \\
\end{srstable}

\resetreqcounter{NvIR}
\begin{srstable}[
  caption = {Требования к интерфейсу навигации.},
]{colspec = {Q[c, m] X[j, m]}}
  Идентификатор & Описание \\
  \reqlabel[ref:nvi1]{NvIR} & Система \textit{Должна} обеспечивать появление элемента управления «Наверх» при осуществлении прокрутки страницы вниз на расстояние, превышающее 300 (триста) пикселей. \\
  \reqlabel[ref:nvi2]{NvIR} & Система \textit{Должна} фиксировать элемент управления «Наверх» в левом нижнем углу области отображения страницы. \\
\end{srstable}

\resetreqcounter{AdIR}
\begin{srstable}[
  caption = {Требования к интерфейсу блока «Совет для вас».},
]{colspec = {Q[c, m] X[j, m]}}
  Идентификатор & Описание \\
  \reqlabel[ref:adi1]{AdIR} & Система \textit{Должна} выполнять асинхронный запрос для загрузки содержимого блока с использованием метода HTTP POST по адресу \texttt{/ajax.php} с параметром \texttt{sovet\_dnya: true}. \\
  \reqlabel[ref:adi2]{AdIR} & Система \textit{Должна} формировать ответ сервера в виде HTML-блока, содержащего изображение карты Таро, наименование карты и текст совета. \\
\end{srstable}

\resetreqcounter{ItIR}
\begin{srstable}[
  caption = {Требования к интерфейсу блока «Это интересно».},
]{colspec = {Q[c, m] X[j, m]}}
  Идентификатор & Описание \\
  \reqlabel[ref:iti1]{ItIR} & Система \textit{Должна} обеспечивать встраивание текста интересного факта в HTML-код страницы при её первой загрузке (технология Server-Side Rendering, SSR). \\
  \reqlabel[ref:iti2]{ItIR} & Система \textit{Должна} выполнять асинхронный запрос к серверу для получения нового интересного факта при активации пользователем кнопки обновления содержимого блока. \\
\end{srstable}

\resetreqcounter{FvIR}
\begin{srstable}[
  caption = {Требования к интерфейсу избранного.},
]{colspec = {Q[c, m] X[j, m]}}
  Идентификатор & Описание \\
  \reqlabel[ref:fvi1]{FvIR} & Система \textit{Должна} осуществлять добавление страницы в список избранного посредством асинхронного запроса с использованием метода HTTP POST по адресу \texttt{/favorite.php}. \\
  \reqlabel[ref:fvi2]{FvIR} & Система \textit{Должна} передавать в составе запроса идентификатор страницы и её наименование. \\
  \reqlabel[ref:fvi3]{FvIR} & Система \textit{Должна} заменять элемент управления, инициирующий добавление в избранное, на текстовую надпись «В избранных» при успешном выполнении операции добавления. \\
  \reqlabel[ref:fvi4]{FvIR} & Система \textit{Должна} обновлять отображаемое значение счётчика избранных страниц при успешном добавлении страницы в избранное. \\
  \reqlabel[ref:fvi5]{FvIR} & Система \textit{Должна} осуществлять переход на страницу \texttt{/izbrannoe} при активации пользователем элемента управления, отображающего счётчик избранного. \\
\end{srstable}

\resetreqcounter{NtIR}
\begin{srstable}[
  caption = {Требования к интерфейсу натальной карты.},
]{colspec = {Q[c, m] X[j, m]}}
  Идентификатор & Описание \\
  \reqlabel[ref:nti1]{NtIR} & Система \textit{Должна} предоставлять для заполнения поля «дата рождения» графический календарь, обеспечивающий возможность выбора года, месяца и дня. \\
  \reqlabel[ref:nti2]{NtIR} & Система \textit{Должна} предоставлять для заполнения поля «время рождения» два выпадающих списка: для выбора часа (в диапазоне от 0 до 23) и для выбора минуты (в диапазоне от 0 до 59). \\
  \reqlabel[ref:nti3]{NtIR} & Система \textit{Должна} обеспечивать функциональность автодополнения для поля «город рождения»: при вводе первых трёх букв отображаются варианты городов, соответствующие введённой последовательности. \\
  \reqlabel[ref:nti4]{NtIR} & Система \textit{Должна} отображать результат расчёта натальной карты асинхронно, на той же странице, без её перезагрузки. \\
\end{srstable}

\resetreqcounter{OmIR}
\begin{srstable}[
  caption = {Требования к интерфейсу раздела «Народные приметы».},
]{colspec = {Q[c, m] X[j, m]}}
  Идентификатор & Описание \\
  \reqlabel[ref:omi1]{OmIR} & Система \textit{Должна} предоставлять возможность выбора конкретной даты посредством гиперссылок на странице раздела «Народные приметы». \\
  \reqlabel[ref:omi2]{OmIR} & Система \textit{Должна} отображать отдельную страницу с информацией, соответствующей выбранной пользователем дате. \\
\end{srstable}

\resetreqcounter{ArIR}
\begin{srstable}[
  caption = {Требования к интерфейсу раздела статей.},
]{colspec = {Q[c, m] X[j, m]}}
  Идентификатор & Описание \\
  \reqlabel[ref:ari1]{ArIR} & Система \textit{Должна} реализовывать выпадающий список оглавления статьи как кликабельный элемент, обеспечивающий навигацию по разделам текста. \\
  \reqlabel[ref:ari2]{ArIR} & Система \textit{Должна} отображать имя и фамилию автора в конце текста каждой статьи. \\
\end{srstable}